\documentclass{article}
\usepackage{enumitem}
\usepackage{hyperref}
\author{Tablebri}

\begin{document}

\title{Documentation for the project BrickGame v1.0 aka Tetris}
\date{}
\maketitle

\section{Introduction}

The BrickGame v1.0 aka Tetris project is an implementation of the classic arcade game Tetris using the C programming language and the ncurses library for the terminal interface.

\section{Project structure}

\subsection{Library Tetris (\texttt{src/brick\_game/tetris})}

\begin{itemize}[label=--]
    \item Source code files that implement the logic of the Tetris game.
    \item Basic functions for interacting with the playing field, manipulating pieces, and handling user input.
    \item Implementation of a finite state machine to formalize the game logic.
\end{itemize}

\subsection{Terminal interface (\texttt{src/gui/cli})}

\begin{itemize}[label=--]
    \item Source code files responsible for rendering the game in the terminal using the ncurses library.
    \item Implementation of rendering the playing field, managing user input and displaying the current state of the game.
\end{itemize}

\section{Build the project}

The project uses a build system \texttt{make} whith Makefile, including the following goals:

\begin{itemize}
    \item \texttt{all}: Build the project.
    \item \texttt{install}: Installing the program on the system.
    \item \texttt{uninstall}: Removing a program from the system.
    \item \texttt{clean}: Cleaning temporary files and folders.
    \item \texttt{dvi}: Creating a file DVI.
    \item \texttt{dist}: Creating an archive containing the necessary files to build and use the program.
\end{itemize}

\section{Runtime requirements}

The project involves the use of the C11 programming language, the gcc compiler and the ncurses library for the terminal interface.

\section{Installation and Startup Instructions}

\begin{enumerate}
    \item \textbf{Installing dependencies:}
        \begin{itemize}
            \item Make sure you have gcc compiler installed.
            \item Install the ncurses library.
        \end{itemize}
    \item \textbf{Build the project:}
        \begin{itemize}
            \item Execute \texttt{make all} to build the project.
        \end{itemize}
    \item \textbf{:}
        \begin{itemize}
            \item Execute \texttt{make install} to install the program on the system.
        \end{itemize}
    \item \textbf{Start:}
        \begin{itemize}
            \item Execute \texttt{make run} to run the program.
        \end{itemize}
\end{enumerate}

\section{Using the program}

\begin{enumerate}
    \item \textbf{Control:}
        \begin{itemize}
            \item Use the left and right arrows to move the shape horizontally.
            \item Press the down key to drop the figure.
            \item Other keys for additional actions.
            \item To rotate the shape, use z or the up arrow.
            \item To pause the game, use p.
            \item To unpause the game, use p again.
            \item To exit the game use q.
        \end{itemize}
    \item \textbf{Game mechanics:}
        \begin{itemize}
            \item Rotating and moving figures.
            \item Acceleration of the figure's fall.
            \item Show the next figure.
            \item Destroying filled lines.
            \item The game ends when the upper limit of the playing field is reached.
        \end{itemize}
    \item \textbf{Ending the game:}
        \begin{itemize}
            \item The game ends when the upper limit of the playing field is reached.
        \end{itemize}
\end{enumerate}

\section{Testing}

The project includes unit tests using the check library. Library test coverage is at least 80\%.

\end{document}
